\documentclass[10pt]{article}
\usepackage[T1]{fontenc}
\usepackage[utf8]{inputenc}
\usepackage[french]{babel}
\usepackage{amsmath,amssymb,mathrsfs,mathpazo}
\usepackage{geometry}
\geometry{a4paper,margin=2cm}
\usepackage{graphicx,array,tabularx,tabularray,colortbl,multirow,multicol,booktabs,enumitem,pifont,ulem,fancybox,diagbox,cancel}
\usepackage{pstricks,pstricks-add,pst-plot,pst-tree,pst-node,pst-text,pst-func,pst-math,pst-bspline,pst-3dplot,pst-eucl,pst-solides3d}
\let\pscalculate\undefined
\usepackage{pst-bezier}
\usepackage{tikz,pgfplots,tkz-tab}
\usetikzlibrary{arrows,arrows.meta,patterns,calc,positioning,trees,decorations.markings,shapes.geometric,shapes.misc,shapes.arrows,intersections,decorations.pathmorphing,decorations.pathreplacing,backgrounds,shadows,babel}
\pgfplotsset{compat=1.18}
\usepgfplotslibrary{fillbetween}
\usepackage[np]{numprint}
\usepackage{listings,makecell,pas-tableur,twoopt}
\let\barre\undefined
\usepackage[pstricks]{tablvar}
\definecolor{darkgreen}{RGB}{0,100,0}
\providecommand{\red}{\color{red}}
\providecommand{\blue}{\color{blue}}
\providecommand{\green}{\color{green}}
\providecommand{\black}{\color{black}}
\providecommand{\white}{\color{white}}
\providecommand{\gray}{\color{gray}}
\providecommand{\cyan}{\color{cyan}}
\providecommand{\magenta}{\color{magenta}}
\providecommand{\yellow}{\color{yellow}}
\providecommand{\violet}{\color{violet}}
\providecommand{\orange}{\color{orange}}
\providecommand{\euro}{\texteuro}
\providecommand{\EUR}{\texteuro}
\providecommand{\ieme}{\textsuperscript{e}}
\providecommand{\degres}{\ensuremath{^\circ}}
\providecommand{\degree}{\ensuremath{^\circ}}
\providecommand{\vv}[1]{\overrightarrow{#1}}
\providecommand{\Coord}[1]{#1}
\providecommand{\decosix}{}
\providecommand{\starredbullet}{\ensuremath{\bullet}}
\pagestyle{empty}
\setlength{\parindent}{0pt}
\usepackage{adjustbox}\DeclareUnicodeCharacter{2236}{:}\DeclareUnicodeCharacter{2190}{\ensuremath{\leftarrow}}
\providecommand{\R}{\mathbb{R}}
\providecommand{\N}{\mathbb{N}}
\providecommand{\Z}{\mathbb{Z}}
\providecommand{\Q}{\mathbb{Q}}
\providecommand{\D}{\mathbb{D}}
\providecommand{\C}{\mathbb{C}}
\providecommand{\e}{{\mathrm{e}}}
\providecommand{\vect}[1]{\overrightarrow{#1}}
\providecommand{\vectt}[1]{\overrightarrow{\mathrm{#1}}}
\providecommand{\vvt}[1]{\overrightarrow{\mathrm{#1}}}
\providecommand{\np}[2][]{{#2~#1}}
\providecommand{\vv}[1]{{\overrightarrow{#1}}}
\renewcommand{\d}{{\mathrm{d}}}
\providecommand{\degres}{{^\circ}}
\providecommand{\degre}{{^\circ}}
\providecommand{\blue}{\color{blue}}
\providecommand{\red}{\color{red}}
\providecommand{\enskip}{\quad}
\providecommand{\llbracket}{\lbrack\!\lbrack}
\providecommand{\rrbracket}{\rbrack\!\rbrack}
\providecommand{\euro}{{\text{€}}}
\providecommand{\no}{{\text{n°}}}
\providecommand{\interval}[2]{{[#1;#2]}}
\providecommand{\intervFO}[2]{{[#1;#2[}}
\providecommand{\dotfill}{\ldots}

\providecommand{\ds}{\displaystyle}
\providecommand{\textding}[1]{#1}
\providecommand{\Oij}{\(\mathrm{O};\vec{\imath},\vec{\jmath}\)}
\providecommand{\Oijk}{\(\mathrm{O};\vec{\imath},\vec{\jmath},\vec{k}\)}
\providecommand{\Ouv}{\(\mathrm{O};\vec{u},\vec{v}\)}
\providecommand{\intervFF}[2]{\ensuremath{[#1;#2]}}
\renewcommand{\euro}{\ifmmode\text{\texteuro}\else\texteuro\fi}
\DeclareUnicodeCharacter{00B0}{\ifmmode^{\circ}\else\textdegree\fi}
\DeclareUnicodeCharacter{20AC}{\ifmmode\text{\texteuro}\else\texteuro\fi}
% Énoncés : Éducation nationale. Transcription et corrigés : APMEP, Denis Vergès et contributeurs. Découpage et classement : CoopMaths, Sébastien Lozano, Éric Elter et contributeurs. Adaptation de la présentation : Maths Explorer.
% CC BY-SA : https://coopmaths.fr/www/about/
% https://coopmaths.fr/alea/static/e3c/2020/tex/e3c_2020_00_sujet56_3_cor.tex
\begin{document}
\noindent\textit{APMEP / CoopMaths. CC BY-SA. Présentation adaptée par Maths Explorer.}\par\medskip
\begin{center}
\psset{labelFontSize=\scriptstyle,showorigin=false}
\begin{pspicture}(-1,-0.5)(7.5,9.5)
\pspolygon*[linewidth=1pt,linecolor=yellow,opacity=0.3](0,6)(8,0) 
\multido{\n=0+1}{10}{\psline[linewidth=0.75pt,linecolor=lightgray](\n,0)(\n,7.2)}
\multido{\n=0+0.2}{48}{\psline[linewidth=0.35pt,linecolor=lightgray](\n,0)(\n,7.2)}
\multido{\n=0+1}{8}{\psline[linewidth=0.75pt,linecolor=lightgray](0,\n)(9.35,\n)}
\multido{\n=0+0.2}{37}{\psline[linewidth=0.35pt,linecolor=lightgray](0,\n)(9.35,\n)}
\psaxes[linewidth=0.95pt,]{->}(0,0)(0,0)(9.5,7.4)
\psdots[dotstyle=+,dotscale =1.4,dotangle=45](4,3)(0,6)(8,0)
 \uput[ur](4,3){E} \uput[ur](8,0){A} \uput[ur](0,6){B} \uput[ur](2.4,2.2){G}
\psline(0,6)(8,0)
\psset{arrowscale=2}
\psline[linewidth=0.8pt,linecolor=darkgray]{->}(0,0)(0,1)\psline[linewidth=0.8pt,linecolor=darkgray]{->}(0,0)(1,0)
\psline[linewidth=1pt,linecolor=red](0,0)(4,3)
\psline[linewidth=1pt,linecolor=red](0,6)(4,0)
\psline[linewidth=1pt,linestyle=dashed,linecolor=blue](0,2)(2.667,2)
\psline[linewidth=1pt,linestyle=dashed,linecolor=blue](2.667,0)(2.667,2)
\uput[dl](0,0){O}\uput[d](0.5,0){$\vec{\imath}$}\uput[l](0,0.5){$\vec{\jmath}$}
\end{pspicture}
\end{center}

\subsection*{1.}

\paragraph{a.} \( \overrightarrow{OA} \cdot \overrightarrow{OB} = 8 \times 0 + 0 \times 6 = 0 \).

\paragraph{b.} Avec \( E(4 \,;\, 3) \) d'où \( \overrightarrow{OE} \begin{pmatrix} 4 \\ 3 \end{pmatrix} \), on a :
\[
\overrightarrow{OA} \cdot \overrightarrow{OE} = 8 \times 4 + 0 \times 3 = 32.
\]

\subsection*{2.}

\paragraph{a.} Le milieu du côté opposé à \( B \) est celui de \([OA]\). Ses coordonnées sont \( (4 \,;\, 0) \).

Or \( 1{,}5 \times 4 + 0 - 6 = 0 \iff 6 - 6 = 0 \) qui est vrai.

De même pour \( B(0 \,;\, 6) : 1{,}5 \times 0 + 6 - 6 = 0 \) est vraie. Donc \( 1{,}5x + y - 6 = 0 \) est une équation cartésienne de la médiane issue du point \( B \) dans le triangle \( OAB \).

\paragraph{b.} La médiane issue de \( O \) dans le triangle \( OAB \) contient \( O \) et le milieu \( I \) de \([AB]\) ; on a \( I(4 \,;\, 3) \).

Puisque cette médiane contient \( O \), une de ses équations est \( y = \alpha x \).

Donc en utilisant les coordonnées de \( I \) : \( 3 = \alpha 4 \iff \alpha = \dfrac{3}{4} = 0{,}75 \). Une équation de cette médiane est \( y = 0{,}75x \).

\subsection*{3.}

\( G \), étant commun aux trois médianes du triangle \( OAB \), est le point d'intersection des deux médianes issues de \( B \) et de \( O \). Les coordonnées de \( G \) vérifient donc les deux équations :
\[
\begin{cases}
1{,}5x + y - 6 = 0 \\
y = 0{,}75x
\end{cases}
\]
D'où :
\begin{align*}
&1{,}5x + 0{,}75x - 6 = 0 \\
\iff &2{,}25x = 6 \\
\iff &x = \dfrac{8}{3},
\end{align*}
puis :
\[
y = 0{,}75x = 0{,}75 \times \dfrac{8}{3} = 2.
\]

Conclusion : \( G\left(\dfrac{8}{3} \,;\, 2\right) \).

\end{document}
